\scanpage{68}
\margin{\deleted{(\unsure{B4 B64})}}
qui définit une triangulation (orientée) avec pondération par les images inverses des pts $0,1,\infty$ (qui forment \inserted{\deleted{par définition}} une partition de l'ens des sommets $\unsure{\}}$ de la triangulation, défini comme l'ens des pts au dessus de $\{0,1,\infty\}$), et en passant au quotient par la conjugaison complexe -- ce qui donne une surface sans bord, grâce à la condition (27). Dans le cas où on ne pose pas cette condition, l'espace quotient
\[
\text{(28)}\qquad X=X'(\mathbb C)/\tau
\]
est une surface à bord compacte (et munie d'une structure conforme...), et on pourrait étendre \inserted{(de façon ess.t unique)} la définition des cartes \inserted{\deleted{finies}} triangulées pondérées sur des surfaces compactes \unsure{à} bord, au cas des surfaces compactes à bord, de telle façon que (26) devienne une équivalence de \uline{catégories}.
\figure{\begin{tikzpicture}\draw (0,0)--(.28,-.38);\end{tikzpicture}}

\textcircled{4.} On peut introduire, sur nos cartes (orientée ou non), des conditions de la forme
\[
\text{(29)}\qquad \rho_i^{\nu_i}=1
\]
\deleted{\unsure{type}} où $\nu_i\in\unsure{\overline{\mathbb N}^{*}|\mathbb N^{*}}$, qui signifie donc que
\margin{67}
